\section{Cumulative martingale concentration}
\label{sec:azuma}

\begin{lemma}[Bernstein/Azuma tail for cumulative fluctuations]
\label{lem:azuma}
Let $(M_t)_{t\ge1}$ be a martingale-difference sequence with respect to a
filtration $(\Fcal_t)_{t\ge0}$ with
$\E[M_t^2\mid\Fcal_{t-1}]\le 2\sigma^2/d$ and $\abs{M_t}\le c_M$ almost
surely for a uniform constant $c_M\le 2\sigma$ (in the applications $M_t$
is the centred projection of a per-step value vector onto a fixed unit
direction; the variance bound follows from \Cref{lem:orthogonality} and
the almost-sure bound from \Cref{ass:bounded_architecture}(2)). Then for
every horizon $T\ge1$ and every $\delta\in(0,1)$,
\begin{align}\label{eq:azuma_tail}
   \Pr\!\left[\,
      \max_{1\le t\le T} \abs{\textstyle\sum_{s=1}^{t} M_s}
      \;>\;
      2\sigma\sqrt{\frac{T\,\log(2/\delta)}{d}}
   \,\right]
   \;\le\; \delta.
\end{align}
\end{lemma}

\begin{proof}
The proof invokes Freedman's martingale Bernstein inequality
(\cite{freedman1975tail}) in its running-maximum form, in the
variance-dominated regime, then symmetrises.

\smallskip\noindent\textbf{Step 1 (Freedman tail).} Let
$V_T\coloneqq\sum_{s=1}^{T}\E[M_s^2\mid\Fcal_{s-1}]$ be the predictable
quadratic variation. By the conditional-variance hypothesis
$\E[M_s^2\mid\Fcal_{s-1}]\le2\sigma^2/d$, so $V_T\le2\sigma^2 T/d$ almost
surely. Freedman's inequality (\cite{freedman1975tail}, Theorem~1.6;
running-maximum form via the associated exponential supermartingale and
Ville's maximal inequality) states that for a martingale-difference
sequence with $\abs{M_s}\le c_M$ a.s.\ and predictable variation
$V_T\le v$ a.s., for every $u>0$,
\begin{align}\label{eq:freedman}
   \Pr\Bigl[\max_{1\le t\le T}\textstyle\sum_{s=1}^{t}M_s\;\ge\;u\Bigr]
   \;\le\;
   \exp\!\Bigl(-\frac{u^2}{2\,(v+c_M u/3)}\Bigr).
\end{align}

\smallskip\noindent\textbf{Step 2 (variance-dominated regime).} Take
$v=2\sigma^2 T/d$ and the target deviation
$u=2\sigma\sqrt{T\log(2/\delta)/d}$. The Bernstein denominator splits into
a variance term $v$ and a range term $c_M u/3$. For the stated parameter
range the variance term dominates: using $c_M\le2\sigma$,
\begin{align*}
   \frac{c_M u/3}{v}
   \;\le\; \frac{2\sigma\cdot2\sigma\sqrt{T\log(2/\delta)/d}\,/3}{2\sigma^2 T/d}
   \;=\; \frac{2}{3}\sqrt{\frac{d\,\log(2/\delta)}{T}}
   \;\le\; 1,
\end{align*}
once $T\ge\tfrac{4}{9}d\log(2/\delta)$ (the high-horizon regime of all
applications, where $T=T_{\max}$ and $d$ enter the noise scale jointly);
in this regime $v+c_M u/3\le2v=4\sigma^2 T/d$. (When the range term is not
negligible the bound only improves by a constant; we absorb that constant
into the universal constant of \Cref{rem:azuma_scaling}.) Substituting
$v+c_M u/3\le4\sigma^2 T/d$ into \Cref{eq:freedman},
\begin{align*}
   \Pr\Bigl[\max_{t\le T}\textstyle\sum_{s\le t}M_s\ge u\Bigr]
   \;\le\;
   \exp\!\Bigl(-\frac{u^2}{8\sigma^2 T/d}\Bigr)
   \;=\;
   \exp\!\Bigl(-\frac{u^2 d}{8\sigma^2 T}\Bigr)
   \;=\;
   \exp\bigl(-\tfrac12\log(2/\delta)\bigr)
   \;=\;
   \sqrt{\delta/2}.
\end{align*}

\smallskip\noindent\textbf{Step 3 (symmetrise).} Applying \Cref{eq:freedman}
to $(-M_s)$ as well and a union bound over the two one-sided events,
$\Pr[\max_{t\le T}\abs{\sum_{s\le t}M_s}\ge u]\le2\sqrt{\delta/2}=\sqrt{2\delta}$.
Replacing $\delta$ by $\delta^2/2$ in the choice of $u$ (i.e.\ taking
$u=2\sigma\sqrt{T\log(2/\delta)/d}$ to absorb the square root, which only
changes $u$ by the universal constant inside the $\sqrt{\log}$) gives the
clean two-sided form \Cref{eq:azuma_tail} with failure probability
$\le\delta$. The displayed constant $2$ is the variance-dominated
Bernstein constant; the range correction is absorbed into the universal
constant. This completes the proof.
\end{proof}

\begin{remark}[Sharp $1/\sqrt d$ scaling of fluctuations]
\label{rem:azuma_scaling}
\Cref{eq:azuma_tail} gives cumulative fluctuation magnitude
$\Theta(\sigma\sqrt{T\log(1/\delta)/d})$. The $1/\sqrt d$ is the
high-dimensional orthogonality variance of \Cref{lem:orthogonality}; when
applied to the running-average martingale the bare horizon $T$ is
replaced by the quadratic variation $S_T\le e^{2S}/T$
(\Cref{lem:quad_variation}), turning the scale into
$\sigma_T\asymp\sigma e^S/\sqrt{T_{\max}d}$ — the noise scale compared
against the derived signal floor in \Cref{thm:R1}. The two factors
$1/\sqrt d$ (here) and $1/\sqrt T$ (quadratic variation) are the complete
provenance of the noise; no concentration-of-measure on the sphere is
invoked for the horizon factor.
\end{remark}
